% ===== PRÉAMBULE CANONIQUE — à coller VERBATIM en tête de CHAQUE .tex =====
\documentclass[11pt,a4paper]{article}
\usepackage[utf8]{inputenc}\usepackage[T1]{fontenc}\usepackage{lmodern}
\usepackage[french]{babel}
\usepackage{amsmath,amssymb,mathtools}\usepackage{icomma}
\usepackage[version=4]{mhchem}          % \ce{...} : équations & formules
\usepackage{siunitx}\sisetup{locale=FR} % \qty{}{}, \si{}, \ang{}
\usepackage{chemfig}                    % \chemfig{...} : structures développées
\usepackage{xcolor}\usepackage{colortbl}
\usepackage{tikz}\usepackage{pgfplots}\pgfplotsset{compat=1.18}
\usepackage[most]{tcolorbox}\usepackage{enumitem}\usepackage{array,booktabs}
\usepackage[a4paper,margin=2.2cm]{geometry}
\usetikzlibrary{arrows.meta,calc,angles,quotes,positioning,patterns,backgrounds,shapes.geometric}
\definecolor{primary}{RGB}{222,49,99}\definecolor{primarydark}{RGB}{150,22,55}
\definecolor{defcol}{RGB}{0,114,178}\definecolor{methcol}{RGB}{52,92,148}
\definecolor{excol}{RGB}{0,135,90}\definecolor{attncol}{RGB}{200,40,55}
\tcbset{boxbase/.style={enhanced,breakable,boxrule=0.9pt,arc=2.5pt,
  left=3mm,right=3mm,top=2.3mm,bottom=2.3mm,fonttitle=\bfseries,coltitle=white}}
% --- boîtes TUTORIEL (noms EXACTS, ne pas renommer) ---
\newtcolorbox{cle}[1][]{boxbase,colback=defcol!6,colframe=defcol,
  title={\textsc{À retenir}\ifx\relax#1\relax\relax\else\textmd{ : \itshape #1}\fi}}
\newtcolorbox{methode}[1][]{boxbase,colback=methcol!7,colframe=methcol,sharp corners,
  title={\textsc{Méthode}\ifx\relax#1\relax\relax\else\textmd{ --- \itshape #1}\fi}}
\newtcolorbox{exemplebox}[1][]{boxbase,colback=excol!5,colframe=excol,
  title={\textsc{Exemple}\ifx\relax#1\relax\relax\else\textmd{ : \itshape #1}\fi}}
\newtcolorbox{piege}[1][]{boxbase,colback=attncol!6,colframe=attncol,
  title={\textsc{Piège}\ifx\relax#1\relax\relax\else\textmd{ : \itshape #1}\fi}}
\newtcolorbox{remarque}[1][]{boxbase,colback=black!4,colframe=black!45,
  title={\textsc{Remarque}\ifx\relax#1\relax\relax\else\textmd{ : \itshape #1}\fi}}
% --- boîtes EXERCICE / CORRIGÉ (noms EXACTS, auto-comptées) ---
\newtcolorbox[auto counter]{exo}[1][]{boxbase,colback=primary!4,colframe=primary,
  title={\textsc{Exercice}~\thetcbcounter\ifx\relax#1\relax\relax\else\textmd{ --- \itshape #1}\fi}}
\newtcolorbox[auto counter]{corrige}[1][]{boxbase,colback=excol!4,colframe=excol,
  title={\textsc{Corrigé}~\thetcbcounter\ifx\relax#1\relax\relax\else\textmd{ --- \itshape #1}\fi}}
\newcommand{\R}{\mathbb{R}}\newcommand{\N}{\mathbb{N}}\newcommand{\Z}{\mathbb{Z}}\newcommand{\ds}{\displaystyle}
% (un \usepackage{fancyhdr}+en-tête est possible ; il est ignoré côté web)
% =========================================================================
\begin{document}
\sisetup{per-mode=symbol}

\begin{corrige}[prédominance d'un diacide : le carbonate]
\begin{enumerate}[label=\alph*)]
  \item À $\mathrm{pH} = 7{,}4$ : $7{,}4 > 6{,}4$ (donc \ce{HCO3-} domine \ce{H2CO3}) et $7{,}4 < 10{,}4$ (donc
        \ce{HCO3-} domine \ce{CO3^{2-}}). L'espèce majoritaire est \textbf{\ce{HCO3-}}.
  \item À $\mathrm{pH} = 11 > 10{,}4$ : c'est \textbf{\ce{CO3^{2-}}} qui prédomine.
\end{enumerate}
\end{corrige}

\begin{corrige}[Henderson dans le sens inverse]
\begin{enumerate}[label=\alph*)]
  \item $\dfrac{[\ce{HCOO-}]}{[\ce{HCOOH}]} = 10^{\,\mathrm{pH}-\mathrm{p}K_a} = 10^{\,4{,}8-3{,}8} = 10^{1} = \mathbf{10}$.
  \item Comme $\mathrm{pH} > \mathrm{p}K_a$, la forme \textbf{basique} \ce{HCOO-} l'emporte : il y en a \textbf{10
        fois plus} que de forme acide \ce{HCOOH}.
\end{enumerate}
\end{corrige}

\begin{corrige}[prévoir une réaction de bout en bout]
\begin{enumerate}[label=\alph*)]
  \item $\ce{CH3COOH + CN- <=> CH3COO- + HCN}$.
  \item Acide réactif : \ce{CH3COOH} ($\mathrm{p}K_a = 4{,}8$) ; acide formé : \ce{HCN} ($\mathrm{p}K_a = 9{,}2$).
        \[ \Delta\mathrm{p}K_a = 9{,}2 - 4{,}8 = 4{,}4,\qquad K_{\text{équi}} = 10^{4{,}4}. \]
  \item $\Delta\mathrm{p}K_a = 4{,}4 > 2$ : la réaction est \textbf{complète} (l'acide acétique, plus fort,
        cède son proton au cyanure).
\end{enumerate}
\end{corrige}

\begin{corrige}[avec quel acide la base réagit-elle le mieux ?]
Ici l'acide formé est toujours \ce{NH4+} ($\mathrm{p}K_a = 9{,}2$), donc $\Delta\mathrm{p}K_a = 9{,}2 - \mathrm{p}K_a(\ce{HA})$.
\begin{enumerate}[label=\alph*)]
  \item \ce{HNO2} : $9{,}2 - 3{,}3 = 5{,}9$ ;\quad \ce{CH3COOH} : $9{,}2 - 4{,}8 = 4{,}4$ ;\quad
        \ce{HClO} : $9{,}2 - 7{,}5 = 1{,}7$.
  \item La réaction est d'autant plus complète que $\Delta\mathrm{p}K_a$ est grand : c'est avec \textbf{\ce{HNO2}}
        (l'acide le plus fort, $\Delta\mathrm{p}K_a = 5{,}9$) que la réaction est la plus complète. Avec \ce{HClO}
        ($\Delta\mathrm{p}K_a = 1{,}7 < 2$) elle n'est qu'équilibrée.
\end{enumerate}
\end{corrige}

\begin{corrige}[fractions d'un diagramme de distribution]
\begin{enumerate}[label=\alph*)]
  \item À $\mathrm{pH} = \mathrm{p}K_a$, on a $[\ce{H3O+}] = K_a$, donc
        $\alpha_{\ce{A-}} = \dfrac{K_a}{K_a + K_a} = \dfrac12 = \mathbf{50\,\%}$.
  \item À $\mathrm{pH} = \mathrm{p}K_a + 1$ : $[\ce{H3O+}] = 10^{-(\mathrm{p}K_a+1)} = K_a/10$, d'où
        \[ \alpha_{\ce{A-}} = \frac{K_a}{K_a/10 + K_a} = \frac{1}{1/10 + 1} = \frac{1}{1{,}1} \approx 0{,}91
           = \mathbf{91\,\%}. \]
        (Une unité de pH au-dessus du $\mathrm{p}K_a$ suffit pour que la base dépasse $90\,\%$.)
\end{enumerate}
\end{corrige}

\end{document}
